% GENERATED FILE. Edit canonical lessons, then run npm run generate:books.
% canonical-source-hash: fnv1a64:af9219a6edb1af31
% canonical-lessons: ES-C353-hospital, ES-C353-banco, ES-C353-farmacia, ES-C353-biblioteca, ES-C353-museo

\chapter{Five Buildings You May Have To Find}
\label{ch:el-hospital-y-el-banco}
\hlchaptermodality{353}

\begin{chapteropening}
\textbf{By the end of this chapter:} \emph{I can say el hospital, el banco, la farmacia, la biblioteca and el museo, and ask where any one of them is.}

Complete ES-C353-museo: name five buildings a town has, and ask a stranger where one of them is.
\end{chapteropening}

\section[el hospital]{el hospital — the guest house that became a sick house}

\label{lesson:ES-C353-hospital}



\begin{quote}
Say \emph{what is left over}, then \emph{scarce}. (\emph{El resto}; \emph{escaso}.)
\end{quote}

\subsection*{What to know first}
\begin{itemize}
\raggedright
  \item la salud — what you go there to get back.
  \item el dolor — what takes you there.
\end{itemize}

\begin{sounds}
\begin{itemize}
\raggedright
  \item \emph{os-pi-TAL}, three beats with the weight on the last. The \emph{h} says nothing at all. $\to$ pronunciation reference
\end{itemize}
\end{sounds}

\begin{cousinweb}
\textbf{el hospital} — \textquotedblleft{}the hospital.\textquotedblright{}

Latin \textbf{hospes} is one of the strangest nouns Rome had. It meant \emph{host} and it meant \emph{guest}, in one word, because the obligation ran both ways: you fed the traveller today and he fed you next year. Neither of you was the important one.

A \textbf{hospitāle} was a house kept for such people. Monasteries ran them along the pilgrim roads, and the guests who stayed longest were the ones too ill to walk on. Over a few centuries the guest house quietly became a sick house, and the name stayed put.

English took the word three separate times, by three separate roads, and ended up with three buildings where Rome had one. Two of them are \textbf{hospital} and \textbf{hostel}; the third is waiting for you in the next chapter. Add \textbf{hospice} and \textbf{hospitality} and it has the set.

There is a darker cousin. Latin \textbf{hostis} started as the stranger too, but the stranger you did \emph{not} take in — so English \textbf{hostile} and \textbf{hospitable} are two answers to the same knock at the door.
\end{cousinweb}

\begin{grammarlens}[title={What you've built}]
The first of five buildings you can ask a stranger for by name.
\end{grammarlens}

\subsection*{Your turn}
Say these aloud:

\begin{itemize}
\raggedright
  \item \emph{el hospital}
  \item \emph{un hospital}
  \item \emph{¿dónde está el hospital?}
\end{itemize}

\begin{tcolorbox}[breakable,colback=teal!4,colframe=teal!35!black,title={Before you move on}]
What two people did \emph{hospes} mean at the same time? (The host and the guest.) And which English word for a cheap place to sleep is the same Latin word? (Hostel.)
\end{tcolorbox}
\section[el banco]{el banco — the bench the money sat on}

\label{lesson:ES-C353-banco}



\begin{quote}
Say \emph{scarce}, then \emph{the hospital}. (\emph{Escaso}; \emph{el hospital}.)
\end{quote}

\subsection*{What to know first}
\begin{itemize}
\raggedright
  \item el dinero — what it holds for you.
  \item la moneda — the small hard piece of it.
\end{itemize}

\begin{sounds}
\begin{itemize}
\raggedright
  \item \emph{BAN-ko}, two beats with the weight on the first. The \emph{b} is soft between vowels and firm at the start. $\to$ pronunciation reference
\end{itemize}
\end{sounds}

\begin{cousinweb}
\textbf{el banco} — \textquotedblleft{}the bank,\textquotedblright{} and also \textquotedblleft{}the bench.\textquotedblright{}

Both meanings are alive in Spanish at once, and that is not an accident: they are the same object, four hundred years apart.

The word is Germanic — a \textbf{bank}, a bench — and Italian took it as \emph{banco}. In the Italian trading cities a man who changed one city's coins for another's set up a plank on trestles in \emph{la plaza} and worked at it in the open. People stopped saying \textquotedblleft{}the man at the bench\textquotedblright{} and started saying \textquotedblleft{}the bench.\textquotedblright{}

When one of them could not pay, his creditors broke the plank. \textbf{Banca rotta}, a broken bench, is where English gets \textbf{bankrupt} — a word about furniture.

The bench is also the seat: Italian \textbf{banchetto}, a little bench, is English \textbf{banquet}, because you sat on a bench to eat. And plain English \textbf{bench} is the same Germanic word, arriving on its own.
\end{cousinweb}

\begin{grammarlens}[title={What you've built}]
A word you now own twice over — the place that holds your money, and the seat in the square outside it.
\end{grammarlens}

\subsection*{Your turn}
Say these aloud:

\begin{itemize}
\raggedright
  \item \emph{el banco}
  \item \emph{un banco}
  \item \emph{el dinero está en el banco}
\end{itemize}

\begin{tcolorbox}[breakable,colback=teal!4,colframe=teal!35!black,title={Before you move on}]
What was a \emph{banco} before it was a bank? (A bench.) And what did the creditors do to it to make somebody bankrupt? (They broke it.)
\end{tcolorbox}
\section[la farmacia]{la farmacia — the shop of the word that would not choose}

\label{lesson:ES-C353-farmacia}



\begin{quote}
Say \emph{the hospital}, then \emph{the bank}. (\emph{El hospital}; \emph{el banco}.)
\end{quote}

\subsection*{What to know first}
\begin{itemize}
\raggedright
  \item el remedio — what you go in for.
  \item la fiebre — one reason to go.
\end{itemize}

\begin{sounds}
\begin{itemize}
\raggedright
  \item \emph{far-MA-sia}, with the weight on the second beat. The \emph{c} before \emph{i} is soft, and \emph{ia} runs together as one glide. $\to$ pronunciation reference
\end{itemize}
\end{sounds}

\begin{cousinweb}
\textbf{la farmacia} — \textquotedblleft{}the pharmacy,\textquotedblright{} the chemist's shop.

Greek \textbf{phármakon} is a famous problem. It means a drug, and Greek used the one noun for the thing that heals and the thing that kills. Translators have argued about it for two thousand years, because a Greek sentence about a \emph{pharmakon} often refuses to tell you which one is meant.

The Greeks were not being careless. The difference between a medicine and a poison is how much of it you take, and a word that covers both is honest about that.

\textbf{Pharmakeía} was the handling of such things, and it walked into Latin, into French, and into Spanish as \textbf{farmacia}. English kept the same undecided root in \textbf{pharmacy}, \textbf{pharmacist} and \textbf{pharmacology}.

Spanish also kept a short technical twin, \emph{el fármaco}, for the substance itself, so the shop and the drug in it are one word wearing two shapes.
\end{cousinweb}

\begin{grammarlens}[title={What you've built}]
The place to go when the hospital is more than you need.
\end{grammarlens}

\subsection*{Your turn}
Say these aloud:

\begin{itemize}
\raggedright
  \item \emph{la farmacia}
  \item \emph{una farmacia}
  \item \emph{¿dónde está la farmacia?}
\end{itemize}

\begin{tcolorbox}[breakable,colback=teal!4,colframe=teal!35!black,title={Before you move on}]
What two opposite things did \emph{phármakon} mean? (A remedy and a poison.) And what decides which one you have? (How much of it you take.)
\end{tcolorbox}
\section[la biblioteca]{la biblioteca — the book-chest, and the port it is named after}

\label{lesson:ES-C353-biblioteca}



\begin{quote}
Say \emph{the bank}, then \emph{the pharmacy}. (\emph{El banco}; \emph{la farmacia}.)
\end{quote}

\subsection*{What to know first}
\begin{itemize}
\raggedright
  \item el libro — what fills it.
  \item leer — what you go there to do.
\end{itemize}

\begin{sounds}
\begin{itemize}
\raggedright
  \item \emph{bi-blio-TE-ka}, four beats with the weight on the third. The \emph{b} is soft, and \emph{io} runs together. $\to$ pronunciation reference
\end{itemize}
\end{sounds}

\begin{cousinweb}
\textbf{la biblioteca} — \textquotedblleft{}the library.\textquotedblright{}

Two Greek words are bolted together here, and both are worth having.

\textbf{Biblíon} is a book — a small \textbf{bíblos}, which was the papyrus you wrote on. And \emph{biblos} is a place name. \textbf{Byblos} was the Phoenician port that shipped Egyptian papyrus to Greece, and the Greeks named the material after the harbour it came from, the way English named \emph{china} after a country. So every book in the world is quietly named after a town in Lebanon.

\textbf{Th\'{\={e}}kē} is a chest, a case, a thing you put things in, from \textbf{tithénai}, to put. So a \emph{bibliotheke} is a book-chest.

That second half went everywhere. An \textbf{apoth\'{\={e}}kē} was a store-room, and it came down into Spanish twice as \textbf{la bodega} and \textbf{la botica}, into French as \textbf{boutique}, and into English as \textbf{apothecary}. A wine cellar, a chemist's, a shop and a druggist are one Greek chest.

English took the first half instead, in \textbf{Bible} and \textbf{bibliography}.
\end{cousinweb}

\begin{grammarlens}[title={What you've built}]
The building where the books live, and a Greek word for a chest you will meet again in a wine cellar.
\end{grammarlens}

\subsection*{Your turn}
Say these aloud:

\begin{itemize}
\raggedright
  \item \emph{la biblioteca}
  \item \emph{el libro está en la biblioteca}
  \item \emph{¿dónde está la biblioteca?}
\end{itemize}

\begin{tcolorbox}[breakable,colback=teal!4,colframe=teal!35!black,title={Before you move on}]
Which port gave its name to every book? (Byblos.) And what does the \emph{-teca} half mean on its own? (A chest, a place to put things.)
\end{tcolorbox}
\section[el museo]{el museo — the house of the Muses, which was once a research institute}

\label{lesson:ES-C353-museo}



\begin{quote}
Say \emph{the pharmacy}, then \emph{the library}. (\emph{La farmacia}; \emph{la biblioteca}.)
\end{quote}

\subsection*{What to know first}
\begin{itemize}
\raggedright
  \item ver — what you go there to do.
  \item viejo — what most of what is in it is.
\end{itemize}

\begin{sounds}
\begin{itemize}
\raggedright
  \item \emph{mu-SE-o}, three beats with the weight on the second. Each vowel keeps its own beat; nothing runs together. $\to$ pronunciation reference
\end{itemize}
\end{sounds}

\begin{cousinweb}
\textbf{el museo} — \textquotedblleft{}the museum.\textquotedblright{}

A \textbf{mouseîon} is a seat of the \textbf{Moûsai}, the Muses — nine goddesses who between them owned poetry, song, dance, history and astronomy. To a Greek those were one field, and the Muse was who you asked for help before starting.

The most famous \emph{mouseion} was at Alexandria, and it was nothing like a museum. It was a walled institution where scholars lived at royal expense, ate together, argued, and worked — with the great library attached to it. Eratosthenes measured the size of the earth there. It was a research institute, and the glass cases came much later.

\textbf{Mousik\'{\={e}} tékhnē} is \textquotedblleft{}the art of the Muses,\textquotedblright{} and English shortened it to \textbf{music}. So a museum and a song are two rooms in the same house.

And when an English poet begs the \textbf{Muse} for a first line, he is asking one of the nine for exactly what the Alexandrians asked for: something worth writing down.
\end{cousinweb}

\begin{grammarlens}[title={What you've built}]
Five buildings, all of them askable for by name: \emph{el hospital}, \emph{el banco}, \emph{la farmacia}, \emph{la biblioteca}, \emph{el museo}.
\end{grammarlens}

\subsection*{Your turn}
Say these aloud:

\begin{itemize}
\raggedright
  \item \emph{el museo}
  \item \emph{el hospital, el banco, la farmacia, la biblioteca}
  \item \emph{¿dónde está el museo?}
\end{itemize}

\begin{tcolorbox}[breakable,colback=teal!4,colframe=teal!35!black,title={Before you move on}]
What was the Mouseion at Alexandria really for? (Scholars lived and worked there.) And which everyday English word comes off the same Muses? (Music.)
\end{tcolorbox}
